\allowdisplaybreaks
\begin{center}
	\hrule
	\vspace{.4cm}
 \begin{tabular*}{\textwidth}{@{}l@{}|@{\extracolsep{0.6in}}r@{}}%
\parbox{4.25in}{\raggedright{\includegraphics[width=.9\linewidth]{ictp-pwf.pdf}}} &
\parbox[c][]{4in}{{\Large\textbf{Md Akiful Islam Zawad} \par}
                    { Brac University \par}
                    { KHA 224, Progati Sarani, Merul Badda, \par}
                    { Dhaka 1212, Bangladesh \par}
                    { \href{ext.akiful.islam@bracu.ac.bd}{ext.akiful.islam@bracu.ac.bd}} \par}
\end{tabular*}\vspace{.3in}
\definecolor{ceruleanblue}{rgb}{0.16, 0.32, 0.6}
	\LARGE\scshape\textbf{\textcolor{ceruleanblue}{Physics for Bangladesh QFT School}}
\end{center}
\hrule\vspace{.25in}
{\large\textbf{Problem Sheet:}\ \textsc{4} \hspace{\hfill} \large\textbf{Due Date:} \today\\
	\hrule}
 %----------------------------
\paragraph*{Problem Set 4} %\hfill \newline
\\
The Lagrangian for Scalar QED (complex scalar field $\varphi$ interacting with field $A_\mu$) is given by
\begin{align}
    L = -\frac{1}{4} F_{\mu\nu}^2 + (D_\mu \varphi)^* (D^\mu \varphi) - m^2 \varphi^* \varphi, 
\end{align}
where $D_\mu = \partial_\mu + ieA_\mu$ is the usual gauge-covariant derivative.
\begin{itemize}
    \item[(a)] Compute the Interaction Lagrangian.
    \item[(b)] Derive the Feynman rules for the vertices and propagators of the above theory.
\end{itemize}
\bigskip\bigskip\hline\hline\bigskip
\begin{enumerate}
    \item [(a)] Compute the Interaction Lagrangian

The covariant derivative $D_\mu = \partial_\mu + ie A_\mu$ is used to describe how the scalar field $\varphi$ interacts with the photon field $A_\mu$. The kinetic term in the Lagrangian is written as:
\begin{align}
    (D_\mu \varphi)^* (D^\mu \varphi) &= \left( \partial_\mu \varphi^* - ie A_\mu \varphi^* \right) \left( \partial^\mu \varphi + ie A^\mu \varphi \right).
\end{align}
This expression expands into four terms, representing the kinetic energy of the scalar field, the interaction between the scalar field and the photon, and a photon-photon-scalar interaction.

We now expand the expression for $(D_\mu \varphi)^* (D^\mu \varphi)$ to reveal all the terms:
\begin{align}
    (D_\mu \varphi)^* (D^\mu \varphi) &= \left( \partial_\mu \varphi^* \partial^\mu \varphi \right) 
    + ie \left( \varphi^* A_\mu \partial^\mu \varphi \right)
    - ie \left( \partial_\mu \varphi^* A^\mu \varphi \right)
    + e^2 \left( A_\mu A^\mu \varphi^* \varphi \right).
\end{align}
Each of these terms has a specific physical meaning:
\begin{itemize}
    \item $\partial_\mu \varphi^* \partial^\mu \varphi$: This term describes the usual kinetic energy of the complex scalar field $\varphi$.
    \item $ie \varphi^* A_\mu \partial^\mu \varphi$: This is an interaction term between the scalar field $\varphi$ and the photon field $A_\mu$. It represents the emission or absorption of a photon by the scalar field.
    \item $-ie \partial_\mu \varphi^* A^\mu \varphi$: This is another interaction term, similar to the previous one, but involving the complex conjugate field $\varphi^*$.
    \item $e^2 A_\mu A^\mu \varphi^* \varphi$: This term represents the interaction between two photons and the scalar field. It describes processes in which two photons interact with the scalar field simultaneously.
\end{itemize}

Now that we have expanded the covariant derivative terms, we can substitute them into the original Lagrangian. The full expression for the Lagrangian becomes:
\begin{align}
    L &= -\frac{1}{4} F_{\mu\nu}^2 + \partial_\mu \varphi^* \partial^\mu \varphi - m^2 \varphi^* \varphi \nonumber \\
    &\quad + ie \varphi^* A_\mu \partial^\mu \varphi - ie \partial_\mu \varphi^* A^\mu \varphi + e^2 A_\mu A^\mu \varphi^* \varphi.
\end{align}
Here, the first term describes the dynamics of the photon field $A_\mu$, the second and third terms describe the dynamics of the scalar field $\varphi$, and the remaining terms represent the interactions between the scalar field and the photon field.

The interaction Lagrangian consists of the terms that describe the interactions between the scalar field and the photon field. These terms are:
\begin{align}
    L_{\text{int}} &= ie \varphi^* A_\mu \partial^\mu \varphi - ie \partial_\mu \varphi^* A^\mu \varphi + e^2 A_\mu A^\mu \varphi^* \varphi. 
\end{align}
This interaction Lagrangian contains two main types of interactions:
\begin{itemize}
    \item The first two terms describe the interaction between the scalar field and a single photon.
    \item The last term describes the interaction between the scalar field and two photons.
\end{itemize}
\item [(b)] Derive the Feynman Rules for Vertices and Propagators

The propagator for the scalar field $\varphi$ can be derived from the free part of the Lagrangian:
\begin{align}
    L_{\text{free}}(\varphi) = \partial_\mu \varphi^* \partial^\mu \varphi - m^2 \varphi^* \varphi.
\end{align}
This is the standard Lagrangian for a complex scalar field. The propagator for the scalar field is given by:
\begin{align}
    \frac{i}{p^2 - m^2 + i\epsilon},
\end{align}
where $p$ is the 4-momentum of the scalar particle, and $m$ is its mass.

The photon propagator is derived from the free part of the Lagrangian for the photon field:
\begin{align}
    L_{\text{free}}(A_\mu) = -\frac{1}{4} F_{\mu\nu}^2,
\end{align}
which gives the propagator for the photon field:
\begin{align}
    \frac{-ig_{\mu\nu}}{q^2 + i\epsilon},
\end{align}
where $q$ is the 4-momentum of the photon.

The terms $ie \varphi^* A_\mu \partial^\mu \varphi$ and $-ie \partial_\mu \varphi^* A^\mu \varphi$ describe the interaction between the scalar field and a single photon. The Feynman rule for the vertex involving one incoming photon and one incoming/outgoing scalar particle is given by:
\begin{align}
    \text{Vertex factor: } \quad ie (p_1^\mu + p_2^\mu), 
\end{align}
where $p_1^\mu$ and $p_2^\mu$ are the 4-momenta of the incoming and outgoing scalar fields.

The term $e^2 A_\mu A^\mu \varphi^* \varphi$ describes the interaction between two photons and the scalar field. The Feynman rule for this vertex, involving two incoming photons and two scalar fields, is given by:
\begin{align}
    \text{Vertex factor: } \quad 2ie^2 g_{\mu\nu},
\end{align}
where $g_{\mu\nu}$ is the metric tensor for the photon-photon interaction.
\end{enumerate}